\documentclass[11pt,letterpaper]{article}

% ---------- Idioma y codificación ----------
\usepackage[utf8]{inputenc}
\usepackage[T1]{fontenc}
\usepackage[spanish,es-tabla]{babel}

% ---------- Paquetes generales ----------
\usepackage[margin=2.5cm]{geometry}
\usepackage{xcolor}
\usepackage{enumitem}
\usepackage{fancyhdr}
\usepackage{titlesec}
\usepackage{tcolorbox}
\tcbuselibrary{skins,breakable}
\usepackage{listings}
\usepackage{graphicx}
\usepackage{array}
\usepackage{longtable}
\usepackage{hyperref}
\usepackage{xurl}
\usepackage{seqsplit}

% ---------- Colores del curso ----------
\definecolor{brandblue}{HTML}{1F4E78}
\definecolor{brandgreen}{HTML}{1DB954}
\definecolor{lightgray}{HTML}{F4F6F8}
\definecolor{codebg}{HTML}{F5F5F5}

% ---------- Encabezado y pie de página ----------
\pagestyle{fancy}
\fancyhf{}
\renewcommand{\headrulewidth}{0.4pt}
\fancyhead[L]{\small Laboratorio de Simulación 1 --- Analítica de Datos}
\fancyhead[R]{\small Módulo 1: Introducción}
\fancyfoot[C]{\small\thepage}

% ---------- Estilo de secciones ----------
\titleformat{\section}
  {\normalfont\Large\bfseries\color{brandblue}}
  {\thesection}{0.6em}{}
\titleformat{\subsection}
  {\normalfont\large\bfseries\color{brandblue}}
  {\thesubsection}{0.6em}{}

% ---------- Cajas de actividades y notas ----------
\newtcolorbox{actividad}[1][]{
  breakable, enhanced, colback=lightgray, colframe=brandblue,
  boxrule=0.8pt, arc=2pt, left=6pt, right=6pt, top=5pt, bottom=5pt,
  title={\faPencil\ Actividad}, fonttitle=\bfseries\color{white},
  colbacktitle=brandblue, #1
}
\newtcolorbox{nota}[1][]{
  breakable, enhanced, colback=white, colframe=brandgreen,
  boxrule=0.8pt, arc=2pt, left=6pt, right=6pt, top=5pt, bottom=5pt,
  title={Nota}, fonttitle=\bfseries\color{white},
  colbacktitle=brandgreen, #1
}

% ---------- Código / comandos ----------
\lstset{
  basicstyle=\ttfamily\small,
  backgroundcolor=\color{codebg},
  breaklines=true,
  frame=single,
  framerule=0.3pt,
  rulecolor=\color{gray!50},
  columns=fullflexible,
  keepspaces=true,
}

% ---------- Metadatos del PDF ----------
\hypersetup{
  pdftitle={Laboratorio de Simulación 1 - Analítica de Datos},
  pdfauthor={Diego Fabián Collazos Huertas},
  pdfsubject={Curso Analítica de Datos - Módulo 1},
  colorlinks=true,
  linkcolor=brandblue,
  urlcolor=brandblue,
  citecolor=brandblue,
}

% Sustituto simple para el icono de lápiz (evita depender de fontawesome)
\newcommand{\faPencil}{\raisebox{-0.1em}{\textbf{\Large\textbullet}}}

\title{\vspace{-2cm}}
\date{}

\begin{document}

% ============================================================
% PORTADA
% ============================================================
\begin{titlepage}
\centering
\vspace*{1.5cm}

{\Large Universidad Nacional de Colombia --- Sede Manizales}\\[0.2cm]
{\large Facultad de Ingeniería y Arquitectura}\\[0.2cm]
{\large Departamento de Ingeniería Eléctrica, Electrónica y Computación}\\[1.2cm]

\rule{\linewidth}{1pt}\\[0.5cm]
{\huge\bfseries\color{brandblue} Laboratorio de Simulación 1}\\[0.4cm]
{\Large Implementación en Orange Data Mining}\\[0.5cm]
\rule{\linewidth}{1pt}\\[1cm]

{\large \textbf{Curso:} Analítica de Datos (Aprendizaje de máquina) --- Código 4200729}\\[0.3cm]
{\large \textbf{Módulo:} 1 --- Introducción}\\[0.3cm]
{\large \textbf{Docente:} Diego Fabián Collazos Huertas}\\[0.3cm]
{\large \textbf{Modalidad:} trabajo en parejas}\\[1.5cm]

\vfill

{\small Este laboratorio hace parte de la nota de \textbf{Laboratorios de simulación (45\% de la evaluación)}, según el programa de la asignatura (\texttt{PROGRAMA.md}).}

\end{titlepage}

% ============================================================
% OBJETIVOS
% ============================================================
\section{Objetivos}

\begin{itemize}[leftmargin=1.2cm]
  \item Reproducir en \textbf{Orange Data Mining} --- mediante programación visual, sin escribir código --- los mismos análisis que se implementaron en Python en los notebooks \texttt{\seqsplit{1.4\_Bases\_de\_Datos\_y\_Datasets\_Ejemplos.ipynb}} y \texttt{\seqsplit{1.6\_Tipos\_de\_Variables\_y\_Visualizacion\_Ejemplos.ipynb}}.
  \item Practicar la carga de datos en distintos formatos (CSV, y datos ya preprocesados desde una base de datos SQL) usando los widgets de Orange.
  \item Configurar una base de datos SQL y conectarla a Orange mediante el widget \textit{SQL Table}.
  \item Identificar y corregir manualmente el tipo de variable de las columnas de un dataset dentro de Orange.
  \item Elegir y construir la visualización adecuada (gráfico de barras vs. histograma) según el tipo de cada variable.
\end{itemize}

\begin{nota}
Este taller \textbf{no reemplaza} los notebooks de Python; los complementa. La idea es que compares, para el mismo problema, qué tan distinto es resolverlo con código (Pandas/Matplotlib) frente a resolverlo con programación visual (Orange) --- ver \texttt{1.2\_Orange\_Data\_Mining.md} para instalar la herramienta si aún no lo has hecho.
\end{nota}

% ============================================================
% INSTRUCCIONES GENERALES
% ============================================================
\section{Instrucciones generales}

\begin{enumerate}[leftmargin=1.2cm]
  \item \textbf{Este laboratorio se desarrolla en parejas} (2 estudiantes), asignadas o conformadas según indique el docente.
  \item Para cada parte del laboratorio, arma el flujo (\textit{workflow}) solicitado en el lienzo de Orange.
  \item Cada vez que veas el ícono \faPencil\ \textbf{Actividad}, documenta la respuesta \textbf{directamente sobre el lienzo de Orange}, usando la herramienta de texto (clic derecho sobre un espacio vacío del lienzo $\to$ \textit{Text}, o el ícono ``T'' de la barra de herramientas) para agregar un cuadro de texto junto a los widgets correspondientes. No se entregan respuestas en un documento aparte: el workflow comentado \textbf{es} la respuesta.
  \item Además, toma una captura de pantalla del resultado de cada actividad (el widget correspondiente ya abierto, mostrando el gráfico o la tabla) para incluirla en el PDF de evidencias.
  \item Consulta la sección \textbf{Entregables} (al final de este documento) para el detalle de qué debes entregar y cómo.
\end{enumerate}

% ============================================================
% PARTE 1
% ============================================================
\section{Parte 1 --- Bases de datos y datasets en Orange}

Esta parte reproduce en Orange lo desarrollado en el notebook \texttt{\seqsplit{1.4\_Bases\_de\_Datos\_y\_Datasets\_Ejemplos.ipynb}}: cargar el mismo dataset en dos formatos distintos (CSV y datos provenientes de una base de datos SQL) y visualizarlo.

\subsection{1.1 Dataset en formato CSV --- \textit{Wine Quality}}

\begin{enumerate}[leftmargin=1.2cm]
  \item Descarga el archivo CSV del repositorio UCI (el mismo que se usó en el notebook):\\
  \url{https://archive.ics.uci.edu/ml/machine-learning-databases/wine-quality/winequality-red.csv}
  \item Abre Orange y crea un workflow nuevo.
  \item Arrastra el widget \textbf{File} al lienzo y ábrelo. Selecciona el archivo descargado.
  \item \textbf{Importante:} este CSV usa punto y coma (\texttt{;}) como separador, no coma. Si Orange no separa las columnas correctamente, revisa la opción de delimitador en el widget \textit{File} (o renombra el archivo con extensión \texttt{.csv} después de verificar el separador; Orange normalmente lo detecta automáticamente).
  \item Conecta el widget \textbf{Data Table} a \textit{File} para verificar que las 1599 filas y 12 columnas se cargaron correctamente.
  \item Agrega el widget \textbf{Distributions}, conéctalo a \textit{File}, y selecciona la variable \texttt{quality}.
\end{enumerate}

\begin{actividad}
Compara el histograma/gráfico de \texttt{quality} que obtuviste en Orange con el resultado de \texttt{vinos['quality'].value\_counts()} del notebook. En un cuadro de texto junto al widget \textit{Distributions}, escribe si coinciden los conteos por nivel de calidad. Incluye en tu PDF de evidencias una captura de este widget.
\end{actividad}

\subsection{1.2 Dataset en formato Raw Text --- \textit{Auto MPG}}

El archivo \texttt{auto-mpg.data} de la UCI viene en \textbf{texto plano sin delimitador fijo} (separado por una cantidad variable de espacios, con \texttt{?} para los valores faltantes y el nombre del carro entre comillas). Orange está diseñado para trabajar con datos ya tabulares, así que --- igual que harías en un proyecto real --- usamos Python/Pandas para la limpieza de este tipo de formato, y Orange para la exploración y visualización posterior:

\begin{enumerate}[leftmargin=1.2cm]
  \item En el notebook \texttt{1.4}, ejecuta hasta la celda que genera \texttt{autos\_limpio} (sección 2) y expórtalo a CSV agregando esta línea:
  \begin{lstlisting}[language=Python]
autos_limpio.to_csv('autos_limpio.csv', index=False)
  \end{lstlisting}
  \item Carga \texttt{autos\_limpio.csv} en Orange con el widget \textbf{File}, igual que en el punto 1.1.
  \item Agrega un widget \textbf{Scatter Plot}: en el eje X selecciona \texttt{caballos\_fuerza}, en el eje Y selecciona \texttt{mpg}, y en \textit{Color} selecciona \texttt{cilindros}.
\end{enumerate}

\begin{actividad}
Compara el \textit{Scatter Plot} de Orange con el gráfico equivalente hecho con Matplotlib en el notebook (potencia vs. eficiencia, coloreado por cilindros). En un cuadro de texto junto al \textit{Scatter Plot}, escribe si se observa la misma tendencia (a mayor potencia, menor eficiencia). Incluye en tu PDF de evidencias una captura de este widget.
\end{actividad}

\begin{nota}
Este es un buen ejemplo de cuándo usar cada herramienta: Python/Pandas es más flexible para \textbf{parsear} formatos de texto no estandarizados; Orange es más ágil para \textbf{explorar visualmente} datos que ya están en formato tabular.
\end{nota}

\subsection{1.3 Base de datos SQL}

A diferencia de Pandas (que puede leer directamente un archivo SQLite con el módulo \texttt{sqlite3}, como se hizo en el notebook), el widget nativo de Orange para bases de datos --- \textbf{SQL Table} --- está diseñado para conectarse a un \textbf{motor cliente-servidor} como PostgreSQL, no a un archivo SQLite local. Por eso, para esta parte del taller vamos a montar la misma información en \textbf{PostgreSQL}.

\subsubsection*{Paso 1 --- Instalar PostgreSQL}

\begin{itemize}[leftmargin=1.2cm]
  \item \textbf{Windows:} descarga el instalador desde \url{https://www.postgresql.org/download/windows/} y sigue el asistente (recuerda la contraseña que le asignes al usuario \texttt{postgres}).
  \item \textbf{macOS:} \texttt{brew install postgresql@16} y luego \texttt{brew services start postgresql@16}.
  \item \textbf{Linux (Ubuntu/Debian):} \texttt{sudo apt install postgresql} y luego \texttt{sudo systemctl start postgresql}.
\end{itemize}

\subsubsection*{Paso 2 --- Crear la base de datos y cargar los datos}

Abre una terminal y conéctate con el cliente \texttt{psql} (te pedirá la contraseña que configuraste):

\begin{lstlisting}[language=bash]
psql -U postgres
\end{lstlisting}

Dentro de \texttt{psql}, crea la base de datos y la tabla, y sal:

\begin{lstlisting}[language=SQL]
CREATE DATABASE curso_analitica;
\c curso_analitica

CREATE TABLE vinos (
    fixed_acidity        NUMERIC,
    volatile_acidity     NUMERIC,
    citric_acid          NUMERIC,
    residual_sugar       NUMERIC,
    chlorides             NUMERIC,
    free_sulfur_dioxide   NUMERIC,
    total_sulfur_dioxide  NUMERIC,
    density                NUMERIC,
    pH                     NUMERIC,
    sulphates              NUMERIC,
    alcohol                NUMERIC,
    quality                INTEGER
);

\q
\end{lstlisting}

Desde la terminal (no dentro de \texttt{psql}), carga el archivo CSV descargado en el punto 1.1 dentro de la tabla que acabas de crear:

\begin{lstlisting}[language=bash]
psql -U postgres -d curso_analitica -c "\copy vinos FROM 'winequality-red.csv' WITH (FORMAT csv, HEADER true, DELIMITER ';')"
\end{lstlisting}

\begin{nota}
Ajusta la ruta del archivo \texttt{winequality-red.csv} según dónde lo hayas guardado. Si \texttt{psql} no está en el \texttt{PATH} de tu sistema, en Windows normalmente lo encuentras en \texttt{C:\textbackslash Program Files\textbackslash PostgreSQL\textbackslash 16\textbackslash bin\textbackslash psql.exe}.
\end{nota}

\subsubsection*{Paso 3 --- Conectar Orange a la base de datos}

\begin{enumerate}[leftmargin=1.2cm]
  \item En Orange, arrastra al lienzo el widget \textbf{SQL Table} (categoría \textit{Data}) y ábrelo.
  \item Completa los campos de conexión:
  \begin{itemize}
    \item \textbf{Type:} PostgreSQL
    \item \textbf{Server:} \texttt{localhost}
    \item \textbf{Database:} \texttt{curso\_analitica}
    \item \textbf{User name} / \textbf{Password:} las de tu usuario \texttt{postgres}
  \end{itemize}
  \item Una vez conectado, en \textbf{Table} selecciona \texttt{vinos} (o escribe una consulta SQL propia en el campo \textit{Custom SQL}, por ejemplo la misma consulta \texttt{GROUP BY} del notebook).
  \item Conecta la salida a un widget \textbf{Data Table} para confirmar que los 1599 registros llegaron correctamente.
\end{enumerate}

\begin{actividad}
En el campo \textit{Custom SQL} del widget \textit{SQL Table}, escribe la misma consulta que usamos en el notebook:
\begin{lstlisting}[language=SQL]
SELECT quality, COUNT(*) AS num_vinos, ROUND(AVG(ph), 2) AS ph_promedio
FROM vinos
GROUP BY quality
ORDER BY quality;
\end{lstlisting}
Conecta el resultado a un widget \textbf{Bar Chart} (o \textbf{Distributions}) mostrando \texttt{num\_vinos} por \texttt{quality}. En un cuadro de texto junto a ese widget, escribe si el resultado coincide con la tabla \texttt{resumen\_sql} del notebook. Incluye en tu PDF de evidencias una captura del workflow completo (SQL Table $\to$ Data Table $\to$ Bar Chart).
\end{actividad}

% ============================================================
% PARTE 2
% ============================================================
\section{Parte 2 --- Tipos de variables y visualización en Orange}

Esta parte reproduce en Orange lo desarrollado en el notebook \texttt{\seqsplit{1.6\_Tipos\_de\_Variables\_y\_Visualizacion\_Ejemplos.ipynb}}: identificar el tipo de cada variable de un dataset y elegir la visualización correcta.

\subsection{2.1 Cargar el dataset de Kaggle}

\begin{enumerate}[leftmargin=1.2cm]
  \item Si aún no lo tienes, descarga el \textit{Spotify Tracks Dataset} desde Kaggle (\url{https://www.kaggle.com/datasets/maharshipandya/-spotify-tracks-dataset}), o reutiliza el archivo \texttt{.csv} que ya descargaste con \texttt{kagglehub} al ejecutar el notebook \texttt{1.6} (queda guardado en la ruta que imprime la celda de descarga).
  \item Carga ese archivo en Orange con el widget \textbf{File}.
\end{enumerate}

\subsection{2.2 Identificación de tipos de variable}

El widget \textbf{File} de Orange detecta automáticamente el tipo de cada columna y lo muestra con un ícono en la vista previa: \textbf{numérico} (círculo), \textbf{categórico} (cuadrado), \textbf{texto} o \textbf{fecha/hora}. Puedes hacer clic sobre el ícono de cada columna para cambiar manualmente su tipo.

\begin{actividad}
Abre la vista previa de columnas del widget \textit{File} y responde, en un cuadro de texto junto al widget:
\begin{enumerate}[label=\alph*)]
  \item ¿Qué tipo detectó Orange para la columna \texttt{track\_genre}? ¿Coincide con la clasificación que hicimos en el notebook (cualitativa nominal)?
  \item ¿Qué tipo detectó Orange para las columnas \texttt{key} y \texttt{mode}? Como vimos en el notebook, aunque se almacenan como números, en realidad son variables \textbf{cualitativas nominales} (una nota musical, y modo mayor/menor). Si Orange las detectó como numéricas, cámbialas manualmente a \textit{categorical} y explica por qué es importante hacer esta corrección antes de analizarlas.
\end{enumerate}
Incluye en tu PDF de evidencias una captura de la vista previa de columnas del widget \textit{File}, mostrando el ajuste que hiciste.
\end{actividad}

\subsection{2.3 Visualización según el tipo de variable}

\begin{enumerate}[leftmargin=1.2cm]
  \item Conecta el widget \textbf{Distributions} al widget \textit{File}.
  \item Selecciona la variable \texttt{track\_genre}: al ser categórica, Orange debe mostrar automáticamente un \textbf{gráfico de barras} con la frecuencia de cada género.
  \item Ahora selecciona la variable \texttt{tempo}: al ser continua, Orange debe mostrar automáticamente un \textbf{histograma}.
\end{enumerate}

\begin{actividad}
Compara ambos gráficos con los del notebook (barras de género musical, histograma de tempo). En un cuadro de texto junto al widget \textit{Distributions}, escribe si el género más frecuente coincide y si la distribución del tempo tiene una forma parecida. Incluye en tu PDF de evidencias dos capturas del widget \textit{Distributions}: una con \texttt{track\_genre} seleccionada y otra con \texttt{tempo}.
\end{actividad}

% ============================================================
% ENTREGABLES
% ============================================================
\section{Entregables}

La pareja debe subir al repositorio del curso (o entregar por el medio que indique el docente) \textbf{dos tipos de archivo}:

\begin{enumerate}[leftmargin=1.2cm]
  \item \textbf{Los archivos \texttt{.ows} de los workflows armados en Orange} (uno por cada parte del laboratorio, o los que hayas usado), \textbf{comentados con la herramienta de texto de Orange} tal como se indicó en cada actividad. El workflow comentado --- no un documento aparte --- es donde se responden las preguntas de cada actividad.
  \item \textbf{Un único PDF con las capturas de pantalla} de los resultados obtenidos en cada una de las 5 actividades (\faPencil), claramente identificadas (por ejemplo, ``Actividad 1 --- Distributions de \texttt{quality}'').
\end{enumerate}

Nombra los archivos con los apellidos de ambos integrantes de la pareja, por ejemplo: \texttt{\seqsplit{apellido1\_apellido2\_parte1.ows}}, \texttt{\seqsplit{apellido1\_apellido2\_parte2.ows}} y \texttt{\seqsplit{apellido1\_apellido2\_evidencias.pdf}}.

\begin{actividad}[title={\faPencil\ Reflexión final}]
Como cuadro de texto de cierre en el workflow de la \textbf{Parte 2}, escriban en pareja una breve reflexión (5-8 líneas): ¿en qué escenarios preferirían usar Orange sobre Python, y en cuáles preferirían Python sobre Orange, a partir de lo que experimentaron en este laboratorio?
\end{actividad}

\begin{nota}
Este laboratorio corresponde a la primera entrega del componente \textbf{Laboratorios de simulación (45\% de la nota)}. Revisa las fechas de entrega vigentes en \texttt{PROGRAMA.md}.
\end{nota}

\end{document}
